\begin{document}

\title{Modello di Ising: simulazione e analisi dei risultati}
\author{Riccardo Pivi}
\maketitle
\tableofcontents

\section{Introduzione}
Motivazione fisica, sistemi complessi, perché Ising è un modello paradigmatico.

\section{Il modello di Ising}
\subsection{Definizione del modello}
\subsection{Hamiltoniana e osservabili fisiche}
Energia, magnetizzazione, simmetria Z2.

\section{Transizioni di fase e magnetizzazione}
\subsection{Transizione di fase nel limite termodinamico}
\subsection{Comportamento della magnetizzazione}
(Eventuale cenno a $T_c$)

\section{Simulazione Monte Carlo del modello di Ising}
\subsection{Algoritmo di Metropolis}
\subsection{Termalizzazione e autocorrelazione}
\subsection{Osservabili misurate}

\section{Unsupervised learning e transizioni di fase: PCA}
\subsection{Configurazioni di spin come dati}
\subsection{Principal Component Analysis}
\subsection{Interpretazione fisica delle componenti principali}

\section{Analisi dei risultati}
Confronto tra osservabili fisiche e PCA, discussione della transizione.


\end{document}